\documentclass[10pt,a4paper]{amsart}
\usepackage[margin=19mm]{geometry}
\usepackage{amsmath,amssymb,mathtools}
\usepackage[scr=rsfs,cal=euler]{mathalfa}
\usepackage{xfrac}
\usepackage[unicode,hidelinks,bookmarks=false]{hyperref}
\newcommand{\ie}{i.e.\ }
\newcommand{\cf}{cf.\ }
\newcommand{\resp}{respectively\ }
\newcommand{\Subsection}{\subsection}
\providecommand{\nobreakdash}{\nobreak\hbox{-}}
\setlength{\emergencystretch}{2em}
\multlinegap0pt
\usepackage{bm}
\usepackage{bbm}
\usepackage{dsfont}
\usepackage{xspace}
\usepackage{cancel}
\usepackage{mathtools}
\usepackage{amsthm}
\makeatletter
\@ifundefined{theorem}{\newtheorem{theorem}{Theorem}}{}
\@ifundefined{lemma}{\newtheorem{lemma}[theorem]{Lemma}}{}
\makeatother
\newcommand{\gr}{\geqslant}
\newcommand{\ls}{\leqslant}
\title[On the deterministic interior body of random polytopes]{On the deterministic interior body of random polytopes}
\author{Minas Pafis, Natalia Tziotziou}
\date{Source: arXiv:2508.15261v1 (CC BY 4.0)}
\begin{document}
\maketitle

\noindent\emph{Source and licence.} On the deterministic interior body of random polytopes. Version arXiv:2508.15261v1 (CC BY 4.0), distributed under the Creative Commons Attribution 4.0 International licence, as recorded in the arXiv licence metadata for this exact version and shown on the abstract page https://arxiv.org/abs/2508.15261v1. Full rights notice: licenses/arxiv-2508.15261-CC-BY-4.0.txt.\par\smallskip

\noindent\textit{Editorial scope.} Counted unit: the Lemma whose source label is \texttt{lem:exact-vc} in the article (source0.tex lines 546--549, source label \texttt{lem:exact-vc}), delivered together with the article's own complete original proof of it (source0.tex lines 551--591, one \texttt{proof} environment of 41 lines). No mathematics is written, repaired or re-derived: statement and proof are the article's own text line for line. Required context retained uncounted: none. Every definition, display and label of those lines is retained unchanged. Omitted from this package and not redistributed: 1--545, 550--550, 592--1438. Nothing inside a retained excerpt is omitted or altered: every retained body is delivered exactly as the article's lines are written, so no omission receipt of any kind accompanies this package. Citations: the retained excerpts carry no citation command at all, so no bibliography entry of the article is transcribed and no reference is redistributed. Reference adaptation: none was needed and none was made; every reference inside the delivered unit resolves inside it, so no unresolved reference or citation is produced. Printed slips kept literal: no printed word, symbol, hypothesis or number of the counted statement or of its proof was altered. Layout and class: the article's own class is \texttt{article} with packages including srcltx, amsfonts, latexsym, amsmath, amssymb, amsthm, fullpage, dsfont, color, hyperref, setspace; the wrapper is the lane's pinned \texttt{amsart} editorial wrapper at 10pt on A4, which restates the theorem-like environments the retained text needs and adds the article's own macro definitions that the retained text needs (\texttt{\textbackslash gr}, \texttt{\textbackslash ls}), with extra wrapper packages (bm, bbm, dsfont, xspace, cancel, mathtools). The delivered numbering is the unit's own. Assets: no retained span contains a figure environment, a graphics-inclusion command, a PDF-page inclusion, a table or a TikZ picture; no image file is copied into this package, the package declares no asset dependency and no image file is redistributed with it. Licence: version arXiv:2508.15261v1 of this article is distributed under the Creative Commons Attribution 4.0 International licence, as recorded in the arXiv licence metadata for this exact version and shown on the abstract page https://arxiv.org/abs/2508.15261v1. A full rights notice is delivered at licenses/arxiv-2508.15261-CC-BY-4.0.txt. Characters: every retained excerpt is pure ASCII, so the delivered engine, XeLaTeX with its default Latin Modern text font, reports no missing character.
\begin{lemma}\label{lem:exact-vc}Let $C$ be a non-empty compact convex set in $\mathbb{R}^n$.
Consider the family $\mathcal{H}(C)$ of all half-spaces $H^+$ for which $H$ is a supporting hyperplane
of $C$ and $C\subseteq H^-$. Then, the ${\rm VC}$-dimension of $\mathcal{H}(C)$ is equal to $n$.
\end{lemma}

\begin{proof}First we show that ${\rm VC}(\mathcal{H}(C))\gr n$. We may assume that $0\in C$ and set
$m=\max\{h_C(u):u\in S^{n-1}\}$. Note that $m\gr 0$. Then, we define $r=(1+m)\sqrt{n}$ and consider
the vectors $y_i=re_i$, $1\ls i\ls n$, where $\{e_1,\ldots ,e_n\}$ is the standard orthonormal
basis of $\mathbb{R}^n$. For every $\sigma=(\sigma_1,\ldots ,\sigma_n)\in\{-1,1\}^n$ we consider the vector
$u_{\sigma}=\frac{1}{\sqrt{n}}\sigma\in S^{n-1}$. Then, for every $1\ls i\ls n$ we have
$$\langle y_i,u_{\sigma}\rangle =\frac{1}{\sqrt{n}}r\sigma_i=(1+m)\sigma_i,$$
and this implies that
$$\langle y_i,u_{\sigma}\rangle >h_C(u_{\sigma})\;\;\hbox{if}\;\sigma_i=1$$
while
$$\langle y_i,u_{\sigma}\rangle <0\ls h_C(u_{\sigma})\;\;\hbox{if}\;\sigma_i=-1.$$
It follows that $\mathcal{H}(C)$ shatters the set $\{y_1,\ldots ,y_n\}$.

It remains to show that ${\rm VC}(\mathcal{H}(C))\ls n$. We assume that there exist distinct
$x_1,\ldots ,x_{n+1}\in\mathbb{R}^n$ such that $\mathcal{H}(C)$ shatters the set $A=\{x_1,\ldots ,x_{n+1}\}$.
Consider an arbitrary point $x\in C$. Then, $x_1,\ldots ,x_{n+1},x$ are affinely dependent, and hence there
exist $t_1,\ldots ,t_{n+1},t\in\mathbb{R}$, not all of them equal to $0$, such that
$$t+\sum_{i=1}^{n+1}t_i=0\quad\hbox{and}\quad tx+\sum_{i=1}^{n+1}t_ix_i=0.$$
Without loss of generality we assume that $t\gr 0$. By the definition, the elements of $\mathcal{H}(C)$
are the half-spaces of the form
$$H^+(h_C(u),u)=\{y\in\mathbb{R}^n:\langle y,u\rangle \gr h_C(u)\},\quad u\in S^{n-1}.$$
For any $1\ls i\ls n+1$ and $u\in S^{n-1}$ we define
$$g_i(u)=\langle x_i,u\rangle -h_C(u).$$
Therefore, $x_i\in H^+(h_C(u),u)$ if and only if $g_i(u)\gr 0$. Note that, for any $u\in S^{n-1}$,
\begin{equation}\label{eq:exact-1}t(\langle x,u\rangle -h_C(u))+\sum_{i=1}^{n+1}t_ig_i(u)
=\Big\langle tx+\sum_{i=1}^{n+1}t_ix_i,u\Big\rangle -\left(t+\sum_{i=1}^{n+1}t_i\right)\,h_C(u)=0.\end{equation}
Assume that there exists $1\ls s\ls n+1$ such that $t_s>0$. We define $I=\{1\ls i\ls n+1:t_i>0\}$
and $J=\{1,\ldots ,n+1\}\setminus I$. Since $\mathcal{H}(C)$ shatters $A$, we
may find $u\in S^{n-1}$ such that $g_i(u)<0$ for all $i\in I$ and $g_i(u)\gr 0$ for all $i\in J$.
Then, combining the fact that $t(\langle x,u\rangle -h_C(u))\ls 0$ (because $t\gr 0$ and $x\in C$) with the fact
that $t_sg_s(u)<0$ and $t_ig_i(u)\ls 0$ for all other $1\ls i\ls n+1$, we see that
$$t(\langle x,u\rangle -h_C(u))+\sum_{i=1}^{n+1}t_ig_i(u)<0,$$
which is a contradiction by \eqref{eq:exact-1}.

This means that $t_i\ls 0$ for all $1\ls i\ls n+1$, and hence $t=-\sum_{i=1}^{n+1}t_i>0$. Then,
$$x=\sum_{i=1}^{n+1}\left(-\frac{t_i}{t}\right)x_i$$
is a convex combination of $x_1,\ldots ,x_{n+1}$. Since $x\in C$ was arbitrary, we conclude that
$$C\subseteq {\rm conv}(A).$$
This leads to a contradiction. Using again the fact that $\mathcal{H}(C)$ shatters $A$, we can find $H^+\in\mathcal{H}(C)$
such that $A\cap H^+=\varnothing$. But this implies that $C\cap H^+=\varnothing$, which cannot be true
because the boundary of $H^+$ is a supporting hyperplane of $C$.
\end{proof}

\end{document}
